\clearpage\thispagestyle{empty}
\pdfbookmark[0]{Quellen zum Juno-Bericht vom November 1804}{juno180411-sources}
\begin{center}\textbf{Quellen zum Juno-Bericht vom November 1804}\end{center}
\noindent Diese KI-Übertragung in heutiges Deutsch folgt dem Auszug aus Franz Xaver von Zachs Bericht in \href{https://gdz.sub.uni-goettingen.de/id/PPN236007467}{Gauß’ \textit{Werken}, Band~VI} (Göttingen, 1874), S.~256--258. Der ursprüngliche Bericht \textit{Fortgesetzte Nachrichten über den neuen Harding’schen Planeten Juno} erschien als Beitrag XXXVIII der \textit{Monatlichen Correspondenz}, Band~X, November 1804, \href{https://www.digitale-sammlungen.de/en/view/bsb10538602?page=477}{S.~463}--\href{https://www.digitale-sammlungen.de/en/view/bsb10538602?page=485}{471}. Der Auszug beginnt mit dem Vergleichsstern auf \href{https://www.digitale-sammlungen.de/en/view/bsb10538602?page=480}{S.~466}. Er zitiert Carl Friedrich Gauß zu seiner zweiten Bestimmung der Elemente und Heinrich Wilhelm Matthias Olbers zu seiner Hypothese, bei der die von Barnaba Oriani berechneten Verschiebungen der Aphelien zur Sprache kommen. Gauß’ Vergleichstabelle steht auf \href{https://www.digitale-sammlungen.de/en/view/bsb10538602?page=483}{S.~469}. Die Werkausgabe beschreibt die Ephemeride von \href{https://www.digitale-sammlungen.de/en/view/bsb10538602?page=484}{S.~470} in eckigen Klammern.
\par\smallskip\noindent Übertragung und Satz wurden mit KI erstellt. Die Digitalisate der Zeitschrift stammen von der Bayerischen Staatsbibliothek.
\clearpage
